\section{Conclusion}
\label{sec:conclusion}

Whether RLVR outperforms SFT for recovering pruned reasoning models appears to depend on how the model was pruned.
Under diversity-preserving pruning at 20--25\% MLP sparsity on Qwen2.5-family models, the two paradigms produce near-identical accuracy (mean $|\Delta| {=} 0.54$\,pp for the four main criteria on GSM8K at 20\%; max $1.8$\,pp across benchmarks and $\alpha$ variants), yet the pruning criterion creates a 4.4$\times$ peak reward gap.
At 20\% sparsity, SFT warmup recovers the bulk of accuracy (from $<$2\% post-prune to 51--64\%); GRPO exposes structural capacity differences but contributes minimal additional accuracy, a pattern that persists at 300 steps (Appendix~\ref{sec:appendix_extended}).
The non-monotonic $\alpha$ dose-response argues against a layer-distribution confound.
A practical implication is that GRPO reward trajectories and $f_0$ diagnostics can serve as probes for evaluating pruning quality, surfacing structural deficiencies that accuracy alone would miss.
Rather than running full RL recovery to estimate final accuracy (which our results show remains constant across criteria), a short GRPO probe of ${\sim}$100 steps suffices to diagnose whether a pruned model retains the structural diversity needed for downstream RL fine-tuning or task-specific RL alignment.
We discuss broader implications in Appendix~\ref{sec:appendix_discussion}.
